\documentclass[conference]{IEEEtran}
\IEEEoverridecommandlockouts
% The preceding line is only needed to identify funding in the first footnote. If that is unneeded, please comment it out.
%Template version as of 6/27/2024

% --- Packages provided by the IEEE conference template ---
\usepackage{cite}
\usepackage{amsmath,amssymb,amsfonts}
\usepackage{algorithmic}
\usepackage{graphicx}
\usepackage{textcomp}
\usepackage{xcolor}

% --- Additional packages required by this paper's content ---
\usepackage{hyperref}
\usepackage{booktabs}
\usepackage{subcaption}
\usepackage{siunitx}
\usepackage{placeins}
\usepackage{multirow}
\usepackage{array}
\usepackage{threeparttable}

% Figures live in ../images relative to this template directory.
\graphicspath{{../}{../images/}}

\def\BibTeX{{\rm B\kern-.05em{\sc i\kern-.025em b}\kern-.08em
    T\kern-.1667em\lower.7ex\hbox{E}\kern-.125emX}}

% Number sections only down to \subsubsection (prevents a), b), c) on \paragraph)
\setcounter{secnumdepth}{3}

\begin{document}

\title{Humanoid Walking: A Reinforcement Learning Case Study}

\author{\IEEEauthorblockN{1\textsuperscript{st} Francisco Pinto}
\IEEEauthorblockA{\textit{Master in Data Science and Engineering (MECD)} \\
\textit{Faculdade de Engenharia da Universidade do Porto (FEUP)}\\
Porto, Portugal \\
up199504009@edu.fe.up.pt}
\and
\IEEEauthorblockN{2\textsuperscript{nd} João Vieira}
\IEEEauthorblockA{\textit{Master in Data Science and Engineering (MECD)} \\
\textit{Faculdade de Engenharia da Universidade do Porto (FEUP)}\\
Porto, Portugal \\
up202107689@edu.fe.up.pt}
}

\maketitle

\begin{abstract}
Humanoid locomotion remains a challenging reinforcement learning problem due to high-dimensional actuation, intermittent contacts, and strong sensitivity to task and environment design.
While policy-gradient methods such as Proximal Policy Optimization (PPO) can reliably learn flat-ground locomotion, extending these policies to uneven terrain and long-horizon navigation tasks frequently fails in practice.
In this paper, we present an empirical study of implementation-level design choices that govern this transition.

We evaluate four PPO configurations for the MuJoCo Humanoid across tasks of increasing difficulty: flat locomotion, stair climbing, waypoint-based circuit navigation, and circuits that include stair segments.
The configurations differ in reward shaping, observation design for terrain and waypoint perception, and termination handling on uneven terrain.
Across 58 training runs totaling approximately 2.6 billion environment steps, we show that baseline PPO reliably solves flat locomotion but requires explicit task-aligned shaping to acquire stable stair-climbing behavior.
For the most challenging circuit-with-stairs task, fixed world-coordinate termination criteria were found to induce premature failures, preventing learning despite otherwise stable control.
Introducing terrain-relative termination handling enables consistent task completion when combined with navigation-aware reward shaping.

The results highlight termination logic as a critical but often overlooked component of humanoid reinforcement learning and provide practical, implementation-level guidance for scaling locomotion policies beyond flat-ground settings.
\end{abstract}

\begin{IEEEkeywords}
Reinforcement Learning, Humanoid Locomotion, MuJoCo, Stair Climbing, Navigation.
\end{IEEEkeywords}

\section{Introduction}

Reinforcement learning (RL) provides a general framework for sequential decision-making in which an agent learns a policy that maximizes cumulative reward through interaction with an environment~\cite{sutton2018reinforcement}.
Within the Markov Decision Process (MDP) formalism, the agent observes state information, selects actions, and receives scalar rewards that guide policy improvement~\cite{silver2015lectures}.
While RL has demonstrated strong performance on simulated continuous-control benchmarks, humanoid locomotion remains particularly challenging due to high-dimensional actuation, intermittent contacts, and sensitivity to environment and task design.

Recent humanoid systems, such as Boston Dynamics' Atlas~\cite{bostondynamics2024atlas} and Unitree's H1~\cite{unitree2024h1}, exhibit dynamic locomotion behaviors that are often supported by extensive simulation training and large-scale computation.
However, even in simulation, policies that perform well on flat terrain frequently fail when extended to uneven terrain or long-horizon navigation tasks.
In practice, this gap is rarely caused by algorithmic limitations alone, but rather by implementation-level design choices.

In this work, we focus on two such choices that repeatedly emerged as critical failure points when moving beyond flat locomotion: \emph{reward design} and \emph{termination criteria}.
Reward functions that are sufficient for flat walking often provide weak or misleading learning signals on contact-rich terrain such as stairs.
Similarly, termination logic tuned to a fixed ground reference can induce premature episode termination when the agent operates on uneven terrain, effectively preventing learning.

\begin{figure}[h!]
    \centering
    
\includegraphics[width=0.5\linewidth]{images/mujuco_humanoid.png}
    \caption{MuJoCo Humanoid model used across all experimental configurations.}
    \label{fig:humanoid_model}
\end{figure}

This paper presents an empirical study of these issues using the MuJoCo Humanoid~\cite{todorov2012mujoco} and Proximal Policy Optimization (PPO)~\cite{schulman2017ppo} implemented in Stable-Baselines3~\cite{raffin2021sb3}.
We evaluate a sequence of tasks with increasing complexity: flat locomotion, stair climbing, waypoint-based circuit navigation, and circuits that include stair segments.
Rather than proposing a new algorithm, our objective is to identify which \emph{implementation-level design decisions} most strongly affect learning stability and task success in this setting.

Our contributions are empirical and practitioner-oriented:
\begin{itemize}
    \item Two custom MuJoCo environments: a configurable stairs environment and a waypoint-based circuit environment with optional stair segments.
    \item A structured comparison of four PPO configurations that progressively introduce reward shaping, navigation structure, and terrain-aware termination handling.
    \item An empirical analysis based on 58 training runs totaling approximately 2.6B environment steps, culminating in a circuit-with-stairs configuration that achieves high task completion under deterministic evaluation.
\end{itemize}

The remainder of the paper is organized as follows. Section~\ref{sec:mdp} formalizes the MDP framework and defines the state, action, and transition structure common to all tasks. Section~\ref{sec:environments} describes the environments and tasks. Section~\ref{sec:method} details the algorithm, reward configurations, and observation design. Section~\ref{sec:results} presents experimental results. Section~\ref{sec:conclusion} discusses conclusions, limitations, and directions for future work.

For reproducibility, the full implementation---including custom environments, training scripts, and configuration files---is publicly available at \url{https://github.com/ghostcat-pro/rl-humanoid.git}.



%% ====================================================================
%%  SECTION 2: MDP FORMULATION
%% ====================================================================
\section{MDP Formulation}
\label{sec:mdp}

We formalize the humanoid locomotion and navigation tasks as a Markov Decision Process $\mathcal{M} = (\mathcal{S}, \mathcal{A}, P, r, \gamma)$. Table~\ref{tab:notation} summarizes the notation used throughout the paper.

\begin{table}[h!]
\caption{MDP notation and symbols.}
\label{tab:notation}
\centering
\small
\begin{tabular}{@{}cl@{}}
\toprule
\textbf{Symbol} & \textbf{Description} \\
\midrule
$\mathcal{S}$ & State space \\
$\mathcal{A}$ & Action space, $a_t \in \mathbb{R}^{17}$ (joint torques) \\
$P(s' \mid s, a)$ & Transition dynamics (specified in Appendix~\ref{app:transition}) \\
$r(s_t, a_t)$ & Scalar reward signal \\
$\gamma$ & Discount factor ($= 0.99$) \\
$\pi_\theta(a \mid s)$ & Stochastic policy parameterized by $\theta$ \\
$V^{\pi}(s)$ & State value function under $\pi$ \\
$G_t$ & Discounted return from step $t$ \\
$dt$ & Control timestep ($= 0.015\,$s) \\
$z_t$ & Torso height (world frame) \\
$z_{\mathrm{rel},t}$ & Torso height relative to local ground \\
$\psi_t$ & Agent heading (yaw) \\
$d_t$ & Distance to current waypoint \\
$\epsilon_t$ & Heading error to current waypoint \\
$\phi_t, \theta_t$ & Torso roll and pitch angles \\
$v^\star$ & Target locomotion speed (S4) \\
\bottomrule
\end{tabular}
\end{table}


\subsection{State Space}

At each discrete control step $t$, the agent receives an observation $o_t$ that serves as the state representation $s_t$. The base observation is \emph{proprioceptive}: joint positions and velocities, torso pose and velocities, and contact forces, following the standard MuJoCo Humanoid convention~\cite{todorov2012mujoco}. Global $(x,y)$ position is excluded from the observation in all configurations.

For tasks requiring exteroceptive perception, the observation is augmented with task-specific signals (detailed in Section~\ref{sec:obs_design}):
a local terrain height grid for stairs and waypoint-relative navigation features for circuit tasks.

\subsection{Action Space}

The action space is continuous, $a_t \in \mathbb{R}^{17}$, where each component is a joint torque command bounded by $[-0.4, 0.4]$ (Figure~\ref{fig:action_space}). All environments use a MuJoCo physics timestep of $0.003\,$s with a frame skip of $5$, yielding an effective control interval of $dt = 0.015\,$s ($\approx 66.7\,$Hz).

\begin{figure}[h!]
    \centering
    \includegraphics[width=0.5\linewidth]{images/action_space.png}
    \caption{Action space: 17 torque-controlled joints of the MuJoCo Humanoid.}
    \label{fig:action_space}
\end{figure}

\subsection{Return and Objective}

The discounted return from step $t$ is
\begin{equation}
\label{eq:return}
G_t = \sum_{k=0}^{\infty} \gamma^k \, r(s_{t+k}, a_{t+k}),
\end{equation}
and the RL objective is to find policy parameters $\theta^{\star}$ maximizing the expected return:
\begin{equation}
\label{eq:objective}
\theta^{\star} = \arg\max_\theta \; \mathbb{E}_{\pi_\theta}\!\left[ G_0 \right].
\end{equation}
The value function $V^{\pi}(s) = \mathbb{E}_{\pi}\!\left[G_t \mid s_t = s\right]$ serves as a variance-reducing baseline in the actor--critic architecture.

\subsection{Episode Termination}
\label{sec:termination}

Episodes end by either \emph{truncation} (time limit reached) or \emph{termination} (health violation). The time limits are 1000 steps for flat locomotion and stairs, and 5000 steps for circuit tasks.

Health-based termination depends on the torso height. In the baseline (S1), the humanoid is declared unhealthy when $z_t$ falls outside a fixed world-coordinate range $[z_{\min}, z_{\max}]$. On uneven terrain, this induces false terminations when the ground itself is elevated. To address this, configurations S2 and S4 employ a \emph{terrain-relative} health check:
\begin{equation}
\label{eq:terrain_rel_z}
z_{\mathrm{rel},t} = z_t - h(x_t, y_t),
\end{equation}
where $h(x_t, y_t)$ is the terrain height at the agent's planar position. The healthy range $[z_{\min}, z_{\max}]$ is then applied to $z_{\mathrm{rel},t}$ instead of $z_t$.



%% ====================================================================
%%  SECTION 3: ENVIRONMENTS AND TASKS
%% ====================================================================
\section{Environments and Tasks}
\label{sec:environments}

We study four tasks of increasing difficulty, each associated with a specific environment:

\begin{enumerate}
    \item \textbf{Flat locomotion:} stable forward walking on flat terrain (\texttt{Humanoid-v5}).
    \item \textbf{Stair climbing:} ascend a configurable staircase (\texttt{HumanoidStairsConfigurable-v0}).
    \item \textbf{Circuit navigation (flat):} follow a waypoint sequence on flat terrain (\texttt{HumanoidCircuit-v0}).
    \item \textbf{Circuit with stairs:} complete a waypoint course including stair segments (\texttt{HumanoidCircuit-v0} with stairs).
\end{enumerate}

\textbf{Humanoid-v5} is the standard Gymnasium flat-terrain benchmark, used as a control condition.

\textbf{HumanoidStairsConfigurable-v0} is a parameterized environment defined by a flat approach distance, a staircase of $N$ steps (with configurable step height and depth), and an optional top platform. In the experiments reported here, $N \in \{8, \dots, 15\}$ and step height ranges from $0.10$ to $0.20\,$m. A run is considered successful if the agent reaches the top step.

\textbf{HumanoidCircuit-v0} is a waypoint-based navigation environment in which the agent must visit a sequence of planar waypoints in order. The course can be flat or include stair segments between waypoints. Side walls constrain the agent to a corridor, preventing trivial obstacle bypassing.

\subsection{Key Challenges}

Beyond flat locomotion, the tasks introduce several compounding difficulties:
\begin{itemize}
    \item \textbf{Multi-objective coordination:} simultaneous optimization of forward progress, directional control, obstacle traversal, and energy efficiency.
    \item \textbf{Reward sparsity:} circuit navigation introduces sparse milestone rewards (waypoint arrival, circuit completion) atop the dense locomotion signal.
    \item \textbf{Contact discontinuities:} foot impacts on stair edges destabilize gait learning.
    \item \textbf{Turning and heading control:} waypoint sequencing requires accurate heading changes---a non-trivial control problem for a torque-driven biped.
\end{itemize}



%% ====================================================================
%%  SECTION 4: METHOD
%% ====================================================================
\section{Method}
\label{sec:method}

\subsection{Algorithm: Proximal Policy Optimization}
\label{sec:ppo}

We use PPO~\cite{schulman2017ppo} implemented in Stable-Baselines3~\cite{raffin2021sb3}, an on-policy actor--critic method that constrains policy updates via the clipped surrogate objective:
\begin{equation}
\label{eq:ppo}
L^{\mathrm{CLIP}}(\theta) = \hat{\mathbb{E}}_t \!\left[ \min\!\left( \rho_t(\theta)\,\hat{A}_t,\; \mathrm{clip}\!\left(\rho_t(\theta),\, 1{-}\epsilon,\, 1{+}\epsilon\right)\hat{A}_t \right) \right],
\end{equation}
where $\rho_t(\theta) = \pi_\theta(a_t \mid s_t) / \pi_{\theta_{\mathrm{old}}}(a_t \mid s_t)$ is the importance-sampling ratio, $\hat{A}_t$ is the GAE advantage estimate~\cite{schulman2017ppo} with $\lambda = 0.95$, and $\epsilon = 0.2$ is the clipping parameter. Table~\ref{tab:hyperparams} lists all training hyperparameters.

\begin{table}[h!]
\caption{PPO training hyperparameters.}
\label{tab:hyperparams}
\centering
\small
\begin{tabular}{@{}ll@{}}
\toprule
\textbf{Parameter} & \textbf{Value} \\
\midrule
Network architecture & MLP: $256 \times 256$, Tanh \\
Learning rate & $3 \times 10^{-4}$ \\
Batch size & 16{,}384 \\
Epochs per update & 10 \\
Rollout length & 4{,}096 steps $\times$ 8 envs \\
Entropy coefficient & 0.005 \\
Value function coefficient & 0.25 \\
Gradient clipping & 0.5 \\
GAE $\lambda$ & 0.95 \\
Discount $\gamma$ & 0.99 \\
Normalization & VecNormalize (obs \& rewards) \\
Training budget & 10M--100M steps (task-dependent) \\
\bottomrule
\end{tabular}
\end{table}


\subsection{Observation Design}
\label{sec:obs_design}

All configurations share a base proprioceptive observation $o_t^{\mathrm{base}}$. Task-specific augmentations are introduced progressively:

\subsubsection{Terrain Grid (S2, S3, S4)}

A local terrain map with $25$ samples ($5$ x-offsets $\times$ $5$ y-offsets) centered on the agent:
\begin{equation}
\label{eq:terrain_grid}
g_t(i,j) = h(x_t + \delta_i,\, y_t + \delta_j) - z_t, \qquad \delta_i, \delta_j \in \{-0.6, -0.3, 0, 0.3, 0.6\},
\end{equation}
where $(x_t,y_t,z_t)$ are torso world coordinates, $h(x,y)$ is terrain height, and $\delta_i,\delta_j$ (meters) are fixed world-frame sampling offsets. Heights are encoded relative to the agent’s current z.


\subsubsection{Waypoint Features (S3, S4)}

Navigation state is encoded as a 5-dimensional vector:
\begin{equation}
\label{eq:waypoint_signal}
w_t = \Big[\,\Delta x_t,\;\Delta y_t,\;\tfrac{n_t}{N_{\mathrm{wp}}},\;\sin(\epsilon_t),\;\cos(\epsilon_t)\,\Big],
\end{equation}

where $(\Delta x_t, \Delta y_t)$ is the displacement to the active waypoint, $n_t/N_{\mathrm{wp}}$ is normalized fraction of completed waypoints,, and $\epsilon_t=\mathrm{wrap}\!\left(\operatorname{atan2}(\Delta y_t,\Delta x_t)-\psi_t\right)$ is heading error (with $\psi_t$ the agent yaw). $\sin(\epsilon_t),\cos(\epsilon_t)$ avoid angle-wrap discontinuities at $\pm\pi$.


Table~\ref{tab:obs_dims} summarizes the observation dimensionality across configurations.

\begin{table}[h!]
\caption{Observation space dimensions per configuration.}
\label{tab:obs_dims}
\centering
\small
\begin{tabular}{@{}lcccc@{}}
\toprule
\textbf{Component} & \textbf{S1} & \textbf{S2} & \textbf{S3} & \textbf{S4} \\
\midrule
Proprioceptive core (S1-equivalent)\footnotemark & 348 & 348 & 348 & 348 \\
Additional proprioceptive signals & --- & 28 & 28 & 28 \\
Terrain grid $g_t$ & --- & 25 & 25 & 25 \\
Waypoint features $w_t$ & --- & --- & 5 & 5 \\
\midrule
\textbf{Total} & \textbf{348} & \textbf{401} & \textbf{406} & \textbf{406} \\
\bottomrule
\end{tabular}
\end{table}
\footnotetext{S1 uses the default Gymnasium \texttt{Humanoid-v5} base observation (348 dimensions). S2--S4 add 28 proprioceptive channels---including extra worldbody inertia, worldbody COM velocity, worldbody contact-force, and root actuation-force signals---yielding a 376-dimensional base before terrain-grid (+25) and waypoint (+5) features are appended.}

\subsection{Reward Design}
\label{sec:reward_design}
The total reward optimized by PPO consists of a task-specific environment reward $r_t^{\mathrm{env}}$ plus a common shaping wrapper applied to all configurations:
\begin{equation}
\label{eq:total_reward}
r_t^{\mathrm{total}} = r_t^{\mathrm{env}} + r_t^{\mathrm{upright}} - c_t^{\mathrm{effort}} - c_t^{\mathrm{lateral}}.
\end{equation}
The wrapper terms, defined in Table~\ref{tab:wrapper}, encourage upright posture, penalize excessive torque, and discourage lateral drift from the corridor centerline $y = 0$.
\begin{table}[h!]
\caption{Common reward-shaping wrapper terms (Eq.~\ref{eq:total_reward}).}
\label{tab:wrapper}
\centering
\small
\renewcommand{\arraystretch}{1.3}
\begin{tabular}{@{}lll@{}}
\toprule
\textbf{Term} & \textbf{Definition} & \textbf{Weight} \\
\midrule
$r_t^{\mathrm{upright}}$ & $w_u \cdot \max\!\left(0,\; q_{w}^2 - q_{x}^2 - q_{y}^2 + q_{z}^2\right)$ & $w_u = 0.5$ \\[2pt]
$c_t^{\mathrm{effort}}$ & $w_e \cdot \lVert a_t \rVert_2^2$ & $w_e = 0.001$ \\[2pt]
$c_t^{\mathrm{lateral}}$ & $w_\ell \cdot |y_t|$ & $w_\ell = 1.0$ \\
\bottomrule
\end{tabular}
\vspace{2pt}
{\raggedright\footnotesize $q = (q_w, q_x, q_y, q_z)$: torso quaternion;\; $a_t$: action vector;\; $y_t$: lateral position.\par}
\end{table}

As noted in Table~\ref{tab:wrapper}, $q=(q_w,q_x,q_y,q_z)$ denotes the torso quaternion (unit, scalar-first, torso orientation in the world frame), $a_t$ is the action vector (policy output at time $t$, i.e., per-joint control commands), and $y_t$ is the lateral position (signed world-frame $y$-displacement from the corridor centerline, $y=0$).

\subsubsection{S1: Baseline (Humanoid-v5)}

The environment reward is the standard Gymnasium \texttt{Humanoid-v5} decomposition:
\begin{equation}
\label{eq:s1_reward}
r_t^{\mathrm{env}} = r_t^{\mathrm{healthy}} + r_t^{\mathrm{forward}} - c_t^{\mathrm{ctrl}} - c_t^{\mathrm{contact}}.
\end{equation}
Table~\ref{tab:reward_base} defines each component and its default weight.


\begin{table}[h!]
\caption{Baseline reward components (S1, Eq.~\ref{eq:s1_reward}).}
\label{tab:reward_base}
\centering
\small
\renewcommand{\arraystretch}{1.3}

\begin{threeparttable}
\begin{tabular}{@{}llr@{}}
\toprule
\textbf{Component} & \textbf{Definition} & \textbf{Weight} \\
\midrule
$r_t^{\mathrm{healthy}}$ & $w_h \cdot \mathbb{1}\!\left[z_{\min} < z_t < z_{\max}\right]$ & $w_h = 5$ \\[2pt]
$r_t^{\mathrm{forward}}$ & $w_f \cdot \dot{x}_t$ \;\; ($\dot{x}_t = \Delta x / dt$) & $w_f = 1.25$ \\[2pt]
$c_t^{\mathrm{ctrl}}$ & $w_c \cdot \lVert a_t \rVert_2^2$ & $w_c = 0.1$ \\[2pt]
$c_t^{\mathrm{contact}}$ & $w_{\mathrm{cf}} \cdot \mathrm{clamp}\!\left(\lVert F_{\mathrm{ext}} \rVert_2^2\right)$\tnote{a} & $w_{\mathrm{cf}} = 5 \!\times\! 10^{-7}$ \\
\bottomrule
\end{tabular}

\begin{tablenotes}[flushleft]
\footnotesize
\item[a] $F_{\mathrm{ext}}$: concatenated external contact force/torque components over all body segments at time $t$. Here, $\mathrm{clamp}(u)=\min(u,10)$ with $u=\lVert F_{\mathrm{ext}} \rVert_2^2$, i.e., the contact term is capped to avoid rare impact spikes dominating the reward.

\end{tablenotes}
\end{threeparttable}

\end{table}



\subsubsection{S2: Stairs Reward Shaping}

The environment reward extends S1 with climbing-specific terms:
\begin{equation}
\label{eq:s2_reward}
r_t^{\mathrm{env}} = r_t^{\mathrm{forward}} + r_t^{\mathrm{height}} + r_t^{\mathrm{step}} + r_t^{\mathrm{healthy}} - c_t^{\mathrm{ctrl}} - c_t^{\mathrm{contact}},
\end{equation}
where $r_t^{\mathrm{height}} = w_h \cdot \max(0,\, z_t - z_{t-1})$ rewards upward progress and $r_t^{\mathrm{step}} = b_s \cdot \Delta n_t$ is a milestone bonus awarded upon reaching a new stair step ($\Delta n_t \in \{0,1\}$). The healthy reward uses the terrain-relative check (Eq.~\ref{eq:terrain_rel_z}).


\subsubsection{S3: Waypoint Navigation (Flat)}

The environment reward replaces forward velocity with navigation-aligned signals:
\begin{equation}
\label{eq:s3_reward}
r_t^{\mathrm{env}} = r_t^{\mathrm{progress}} + r_t^{\mathrm{dir}} + r_t^{\mathrm{heading}} + r_t^{\mathrm{bonus}} + r_t^{\mathrm{healthy}} - c_t^{\mathrm{ctrl}} - c_t^{\mathrm{contact}}.
\end{equation}
Table~\ref{tab:reward_nav} defines the navigation-specific components.

\begin{table}[h!]
\caption{Navigation reward components (S3, Eq.~\ref{eq:s3_reward}; reused in S4 Table~\ref{tab:s4_reward}).}
\label{tab:reward_nav}
\centering
\small
\renewcommand{\arraystretch}{1.3}
\begin{tabular}{@{}p{1.4cm}p{5.8cm}@{}}
\toprule
\textbf{Term} & \textbf{Definition} \\
\midrule
$r_t^{\mathrm{progress}}$ & $w_p \cdot (d_{t-1} - d_t)$: reduction in waypoint distance. \\[2pt]
$r_t^{\mathrm{dir}}$ & $w_f \cdot \max\!\big(0,\, \dot{\mathbf{x}}^{xy}_t \!\cdot\! \hat{v}_t\big)$: velocity projected onto waypoint direction. Zero when $d_t \leq 0.1\,$m. \\[2pt]
$r_t^{\mathrm{heading}}$ & $w_\psi \cdot \cos(\epsilon_t) \cdot \min(\|\dot{\mathbf{x}}^{xy}_t\|, 2.0)$: heading alignment scaled by speed. Active only when $d_t > 0.1$ and $\|\dot{\mathbf{x}}^{xy}_t\| > 0.1$. \\[2pt]
$r_t^{\mathrm{bonus}}$ & $b_{\mathrm{wp}} \cdot \mathbb{1}[d_t < \tau] + b_{\mathrm{circ}} \cdot \mathbb{1}[\text{all waypoints reached}]$: sparse milestone and completion bonuses. \\
\bottomrule
\end{tabular}
\end{table}

In Table~\ref{tab:reward_nav}, $d_t$ is distance to the active waypoint, $\hat{v}_t$ is the unit direction to it, $\dot{\mathbf{x}}^{xy}_t$ is planar velocity, $\epsilon_t$ is heading error, $\tau$ is waypoint reach threshold, and $w_p,w_f,w_\psi,b_{\mathrm{wp}},b_{\mathrm{circ}}$ are scalar weights. $r_t^{\mathrm{progress}}$ rewards reducing distance, $r_t^{\mathrm{dir}}$ rewards moving toward the waypoint, $r_t^{\mathrm{heading}}$ rewards facing the waypoint (speed-scaled), and $r_t^{\mathrm{bonus}}$ gives waypoint and full-circuit bonuses.



\subsubsection{S4: Circuit with Stairs}

This configuration combines the navigation shaping of S3 with the height-progress term from S2, and introduces two additional regularization terms that proved necessary for stable multi-skill execution: a balance reward and a speed-regulation term. Terrain-relative termination (Eq.~\ref{eq:terrain_rel_z}) is applied.

The final S4 environment reward is
\begin{equation}
\label{eq:s4_reward}
\begin{aligned}
r_t^{\mathrm{env}} ={}& r_t^{\mathrm{progress}} + r_t^{\mathrm{dir}} + r_t^{\mathrm{heading}} + r_t^{\mathrm{balance}} + r_t^{\mathrm{speed}} \\
& + r_t^{\mathrm{height}} + r_t^{\mathrm{bonus}} + r_t^{\mathrm{healthy}} - c_t^{\mathrm{ctrl}} - c_t^{\mathrm{contact}}.
\end{aligned}
\end{equation}
The stability and speed-regularization terms are
\begin{equation}
\label{eq:s4_balance_speed}
\begin{aligned}
r_t^{\mathrm{balance}} &= w_b \cos\!\left(\sqrt{\phi_t^2 + \theta_t^2}\right), \\
r_t^{\mathrm{speed}} &= w_v \exp\!\left(-\frac{1}{2}\left(\frac{\|\dot{\mathbf{x}}_t^{xy}\| - v^\star}{\sigma_v}\right)^2\right),
\end{aligned}
\end{equation}
where $\phi_t$ and $\theta_t$ are the torso roll and pitch extracted from the torso quaternion. The balance term rewards upright posture by penalizing deviations from the vertical axis via the composite tilt angle $\sqrt{\phi_t^2 + \theta_t^2}$. The speed term shapes a Gaussian preference around a target velocity $v^\star$, discouraging both excessively slow locomotion (which stalls navigation progress) and excessively fast movement (which destabilizes stair traversal).

Table~\ref{tab:s4_reward} reports the full S4 specification, including all terms, coefficients, and thresholds used in the final configuration. The total reward optimized by PPO in S4 remains Eq.~\ref{eq:total_reward}, i.e., Eq.~\ref{eq:s4_reward} plus the common wrapper terms from Table~\ref{tab:wrapper} ($w_u = 0.5$, $w_e = 0.001$, $w_\ell = 1.0$).



\subsection{Configuration Summary}

Table~\ref{tab:solutions} summarizes the design choices across all four configurations.

\begin{table}[h!]
\caption{Design choices across configurations S1--S4.}
\label{tab:solutions}
\centering
\small
\begin{tabular}{@{}lcccc@{}}
\toprule
\textbf{Feature} & \textbf{S1} & \textbf{S2} & \textbf{S3} & \textbf{S4} \\
\midrule
Terrain grid observation & --- & \checkmark & \checkmark & \checkmark \\
Height-based climbing term & --- & \checkmark & --- & \checkmark \\
Step milestone bonus & --- & \checkmark & --- & --- \\
Waypoint signals/rewards & --- & --- & \checkmark & \checkmark \\
Terrain-relative termination & --- & \checkmark & --- & \checkmark \\
Stability/speed shaping  & --- & --- & --- & \checkmark \\
\bottomrule
\end{tabular}
\end{table}

\textbf{S1 (Baseline).}
Standard PPO on \texttt{Humanoid-v5} with default rewards and world-coordinate termination. Serves as the flat-locomotion control condition.

\textbf{S2 (Stairs).}
Introduces task-specific climbing rewards, terrain perception, and terrain-relative health checking.

\textbf{S3 (Circuit, flat).}
Extends S2 with waypoint navigation signals and direction-aligned rewards on flat terrain, following a progressive development process from simpler to harder tasks. Results are reported from the corresponding task-specific training runs.

\textbf{S4 (Circuit + stairs).}
Integrates navigation (S3) and climbing (S2) rewards with balance and speed regularization (Eq.~\ref{eq:s4_balance_speed}) and terrain-relative termination. Table~\ref{tab:s4_reward} in Appendix~\ref{app:s4} provides the complete reward specification.


%% ====================================================================
%%  SECTION 5: RESULTS
%% ====================================================================
\section{Results}
\label{sec:results}

We conducted 58 training runs totaling approximately 2.6B environment steps. These runs informed iterative refinements of reward structure, observations, and termination logic. Final results are reported from task-specific training runs, as incremental training across tasks did not reliably preserve previously learned behaviors. Table~\ref{tab:final_runs} identifies the final training runs for each task, and Table~\ref{tab:summary_results} summarizes evaluation performance.

\begin{table}[h!]
\centering
\caption{Final training runs reported in this study. All runs use seed 42.}
\label{tab:final_runs}
\small
\begin{tabular}{@{}llrr@{}}
\toprule
\textbf{Task}  & \textbf{Train Steps} & \textbf{Best Ckpt.} \\
\midrule
S1: Flat                          & 50M & 42M \\
S2: Stairs       & 30M & 29M \\
S3: Circuit              & 80M & 80M \\
S4: Circ.+Stairs         & 80M & 64M \\
\bottomrule
\end{tabular}
\vspace{2pt}

{\raggedright\footnotesize Best Ckpt.: timestep at which the best evaluation checkpoint was selected.\par}
\end{table}

\begin{table}[h!]
\centering
\caption{Summary of evaluation results across all four tasks.}
\label{tab:summary_results}
\small
\begin{tabular}{@{}lrrrr@{}}
\toprule
\textbf{Task} & \textbf{Mean Reward} & \textbf{Std} & \textbf{CV (\%)} & \textbf{Avg.\ Len.} \\
\midrule
Flat locomotion     & 8{,}899   & 406   & 4.6   & 1{,}000 \\
Stair climbing      & 16{,}210  & 226   & 1.4   & 1{,}000 \\
Circuit (flat)      & 7{,}841   & 925   & 11.8  & 524 \\
Circuit + stairs    & 10{,}881  & 105   & 0.96  & 627 \\
\bottomrule
\end{tabular}
\vspace{2pt}

{\raggedright\footnotesize Avg.\ Len.: mean episode length in control steps. CV: coefficient of variation.\par}
\end{table}


\subsection{Flat Locomotion (S1)}

The best flat-locomotion model achieved a mean evaluation reward of $8{,}899 \pm 406$ (CV $= 4.6\%$), with all episodes completing the full 1{,}000-step horizon. Evaluation performance retained approximately 89\% of the peak training reward ($9{,}963$), indicating minimal degradation under deterministic execution.

The training curve (Figure~\ref{fig:s1_curve}) shows rapid initial learning during the first 10M steps, followed by continued improvement up to approximately 40M steps where the mean episode reward plateaus near $8{,}000$--$8{,}500$. The best evaluation checkpoint was selected at 42M timesteps out of a 50M-step budget (Table~\ref{tab:final_runs}), confirming that the policy had effectively converged. Residual oscillations in the plateau region reflect the stochasticity inherent in on-policy rollout collection but do not indicate instability, as the evaluation variance remains low.

\begin{figure}[h!]
    \centering
    \includegraphics[width=0.8\linewidth]{images/S1_Baseline.png}
    \caption{S1 -- Flat locomotion: mean episode reward (\texttt{rollout/ep\_rew\_mean}) over 50M training steps. The policy converges after approximately 40M steps with stable plateau behavior.}
    \label{fig:s1_curve}
\end{figure}


\subsection{Stair Climbing (S2)}

The agent achieved a final evaluation reward of $16{,}210 \pm 226$ (CV $= 1.4\%$), with near-perfect task completion. From 22M timesteps onward, evaluation episodes consistently reached the full 1{,}000-step horizon, corresponding to successful staircase ascent and stable platform traversal. At peak performance (29M timesteps), all episodes completed successfully with reward dispersion below 0.4\%.

The training curve (Figure~\ref{fig:s2_curve}) exhibits the steepest learning trajectory among all configurations, with mean reward rising from approximately $2{,}000$ to $10{,}000$ within the first 10M steps. The policy reaches near-asymptotic performance by 20M steps, and the best checkpoint was obtained at 29M out of a 30M-step budget (Table~\ref{tab:final_runs}). The rapid convergence and low final variance indicate that the combination of height-progress rewards, step milestones, and terrain-relative termination provides a well-shaped learning signal for the stair-climbing task. Notably, this configuration required the shortest training budget (30M steps) among all tasks, despite the added complexity of contact-rich terrain.

\begin{figure}[h!]
    \centering
    \includegraphics[width=0.8\linewidth]{images/S2_Stairs.png}
    \caption{S2 -- Stair climbing: mean episode reward over 30M training steps. Rapid convergence with the best checkpoint at 29M steps reflects the effectiveness of climbing-specific reward shaping.}
    \label{fig:s2_curve}
\end{figure}


\subsection{Circuit Navigation, Flat (S3)}

The flat circuit task yielded a mean reward of $7{,}841 \pm 925$ (CV $= 11.8\%$), reflecting partial mastery. The tests confirmed that full circuit traversal occurred in the majority of rollouts, but not consistently across all episodes. Most trajectories navigated two to three waypoints successfully, with failures arising during later turns. The average episode length of 524 steps indicates sustained locomotion rather than consistent circuit completion. This higher variance highlights the sensitivity of long-horizon waypoint sequencing to turning accuracy and heading alignment.

The training curve (Figure~\ref{fig:s3_curve}) reveals a qualitatively different learning profile from the preceding tasks. The mean reward increases steadily throughout the full 80M-step budget without reaching a clear plateau, suggesting that the policy had not fully converged at training termination. Consistent with this observation, the best evaluation checkpoint coincides with the final timestep (80M; Table~\ref{tab:final_runs}), indicating that additional training could yield further improvement. The gradual, near-linear growth---in contrast to the rapid saturation observed in S1 and S2---reflects the increased difficulty of the long-horizon navigation objective, where the sparse waypoint bonuses provide a weaker learning signal than the dense forward-velocity and height-progress rewards available in flat and stair tasks. However, as this work already encompasses multiple tasks, and ablations under a fixed compute budget, extending training horizons or performing broader hyperparameter exploration for this setting would have imposed a substantial computational and wall-clock cost (approximately 2--3 hours per additional training run), and was therefore deferred to future work.


\begin{figure}[h!]
    \centering
    \includegraphics[width=0.8\linewidth]{images/S3_Circuit.png}
    \caption{S3 -- Circuit navigation (flat): mean episode reward over 80M training steps. The absence of a clear plateau and the best checkpoint at 80M suggest the policy had not fully converged.}
    \label{fig:s3_curve}
\end{figure}


\subsection{Circuit with Stairs (S4)}

The combined task achieved a mean reward of $10{,}881 \pm 105$ (CV $= 0.96\%$), the lowest variance across all tasks. Rewards were tightly clustered between 10{,}750 and 10{,}998, with performance stabilizing after approximately 72M timesteps. The average episode length of 627 steps reflects successful traversal of both waypoint sequences and stair segments.

The training curve (Figure~\ref{fig:s4_curve}) shows rapid early learning followed by a plateau near $7{,}000$--$8{,}000$ from approximately 5M steps onward, with persistent oscillations in the plateau region. The best evaluation checkpoint was selected at 64M out of 80M training steps (Table~\ref{tab:final_runs}), indicating that peak performance was reached before training completion and that later updates did not yield further gains. The oscillatory plateau behavior, more pronounced than in S1, likely reflects the multi-objective nature of the combined task: the policy must simultaneously satisfy navigation, climbing, and stability objectives, and small perturbations in any one component can temporarily degrade the composite reward. Despite this training-time variability, the deterministic evaluation variance is the lowest across all configurations (CV $= 0.96\%$), confirming that the learned policy is robust at deployment time.

This result demonstrates that the terrain-relative termination mechanism (Eq.~\ref{eq:terrain_rel_z}) was essential: without it, the same reward configuration repeatedly failed due to premature episode termination on elevated terrain.

\begin{figure}[h!]
    \centering
    \includegraphics[width=0.8\linewidth]{images/S4_Circuit_Stairs.png}
    \caption{S4 -- Circuit with stairs: mean episode reward over 80M training steps. The best checkpoint at 64M and persistent plateau oscillations reflect the multi-objective complexity of the combined task.}
    \label{fig:s4_curve}
\end{figure}

\subsection{Convergence Analysis}

The training dynamics across the four configurations reveal a clear relationship between task complexity, convergence behavior, and sample efficiency (Figures~\ref{fig:s1_curve}--\ref{fig:s4_curve}). S2 (stairs) converged fastest despite requiring terrain interaction, which we attribute to the dense, well-aligned reward signal provided by height-progress and step-milestone terms. S1 (flat) converged reliably but required more samples due to the relatively weak shaping of the default \texttt{Humanoid-v5} reward. S3 (flat circuit) exhibited the slowest convergence with no clear asymptote, consistent with the sparsity of waypoint-level rewards in a long-horizon task. S4 (circuit with stairs) reached a stable plateau but with higher training-time variance than single-objective tasks, reflecting the tension between competing reward components.

A notable pattern is the discrepancy between training-time oscillations and evaluation-time stability. S4 exhibits the most volatile training curve yet produces the lowest evaluation CV ($0.96\%$), while S3 shows a smoother training trajectory but the highest evaluation CV ($11.8\%$). This suggests that training-curve smoothness is a poor predictor of deployment robustness; instead, the decisive factors are whether the reward structure adequately captures all task-relevant objectives and whether termination logic is consistent with the operating terrain.


%% ====================================================================
%%  SECTION 6: CONCLUSIONS
%% ====================================================================
\section{Conclusions}
\label{sec:conclusion}

This work investigated how implementation-level design choices influence the ability of PPO-trained humanoid agents to scale from flat locomotion to stair climbing and long-horizon waypoint navigation.
Across 58 training runs and approximately 2.6B environment steps, we identified three principal findings.

\textbf{Reward design must be task-aligned.}
Standard forward-velocity rewards are sufficient for flat locomotion but fail on contact-rich terrain. Explicit height-progress terms and step milestones were required for stable stair-climbing behavior. The training dynamics corroborate this: S2 achieved the fastest convergence (best checkpoint at 29M of 30M steps; Figure~\ref{fig:s2_curve}) precisely because its reward components are tightly aligned with the task objective, whereas S3 failed to converge within 80M steps (Figure~\ref{fig:s3_curve}) under sparser waypoint-level signals.

\textbf{Termination logic is a first-class design decision.}
The most critical finding emerged in the circuit-with-stairs task: policies repeatedly failed under fixed world-coordinate termination despite exhibiting otherwise stable locomotion. Terrain-relative termination (Eq.~\ref{eq:terrain_rel_z}) was necessary to enable learning on uneven terrain. With this mechanism in place, S4 achieved the lowest evaluation variance of all configurations (CV $= 0.96\%$; Table~\ref{tab:summary_results}), despite exhibiting the most oscillatory training curve (Figure~\ref{fig:s4_curve}).

\textbf{Naive curriculum transfer is fragile.}
Incremental task training---learning stair climbing first and then extending to circuit navigation---did not reliably preserve previously acquired skills. Without explicit mechanisms to protect earlier competencies, catastrophic forgetting occurred upon introduction of new objectives. The final results were obtained through task-specific training runs (Table~\ref{tab:final_runs}), each trained independently with a common seed.

\textbf{Training-curve smoothness does not predict deployment robustness.}
Comparing Figures~\ref{fig:s3_curve} and~\ref{fig:s4_curve}, S3 exhibits a smoother training trajectory yet produces the highest evaluation variance (CV $= 11.8\%$), while S4's volatile training curve yields the most consistent evaluation performance. This underscores that the quality of reward shaping and termination design---rather than apparent training stability---determines final policy reliability.

\subsection{Limitations}

The final policies rely on privileged terrain observations (local height grids), limiting direct sim-to-real applicability. Evaluation was deterministic over a limited number of episodes for intermediate configurations. All reported runs use a single seed (Table~\ref{tab:final_runs}); multi-seed statistical analysis would strengthen the generality of the conclusions. Generalization to unseen stair geometries and robustness under physical perturbations were not systematically tested. \emph{Moreover, because we iteratively designed rewards and tuned hyperparameters across four distinct tasks, the computational budget was a binding constraint: each training run required approximately 3 hours, which precluded exhaustive per-task sweeps and extended multi-seed experimentation.}


\subsection{Future Work}

Promising directions include reducing reliance on privileged information (e.g., via depth-image policies), evaluating generalization across terrain variations with domain randomization, multi-seed training campaigns for statistical robustness, and incorporating skill-preservation mechanisms such as modular policies or constrained optimization for scalable multi-task humanoid learning.


\bibliographystyle{IEEEtran}
\bibliography{../references}

\appendices

\section{Supplementary Material}
\label{app:supp}

This work includes an alternative PPO training pipeline using TorchRL, which is fully functional and available in the public repository. The motivation was to preserve the option of lower-level intervention in the training stack (e.g., tighter control over rollout collection or loss terms). In the end, Stable-Baselines3 proved sufficient for the reported experiments.

An implementation of Soft Actor-Critic (SAC) was also developed and tested in preliminary runs, but no systematic hyperparameter tuning was performed. The SAC implementation is provided in the repository as a foundation for future work.


\section{Transition Dynamics Specification}
\label{app:transition}

To make transition dynamics explicit, we write
\begin{equation}
    s_{t+1} \sim P_{\theta_{\mathrm{phys}}}(\cdot \mid s_t, a_t),
\end{equation}
where $\theta_{\mathrm{phys}}$ denotes the MuJoCo physics configuration used in each environment.
Because MuJoCo uses soft-contact dynamics, solver and contact settings are part of $P$ and therefore part of the experimental specification.

For S2--S4 (custom environments), the parameters in Table~\ref{tab:transition_params} are fixed in the environment XML.
S1 (\texttt{Humanoid-v5}) uses Gymnasium defaults with no manual physics override in this repository.
Any MuJoCo options not listed in Table~\ref{tab:transition_params} remain at simulator defaults for the MuJoCo version used at runtime.

\begin{table}[h!]
\caption{Physics parameters defining $P_{\theta_{\mathrm{phys}}}$ in S2--S4.}
\label{tab:transition_params}
\centering
\footnotesize
\begin{tabular}{@{}ll@{}}
\toprule
\textbf{Parameter} & \textbf{Value} \\
\midrule
Integrator & \texttt{RK4} \\
Constraint solver & \texttt{PGS} \\
Solver iterations & 50 per physics step \\
Physics timestep & 0.003 s \\
Frame skip & 5 \\
Control interval & 0.015 s \\
Terrain contact dimensionality & \texttt{condim=3} \\
Terrain friction (slide, torsion, roll) & $(1.0, 0.1, 0.1)$ \\
Default geom collision margin & 0.001 \\
Actuator control range & $[-0.4, 0.4]$ \\
Unspecified contact/solver options & MuJoCo defaults (\texttt{solref}, \texttt{solimp}, tolerances, etc.) \\
\bottomrule
\end{tabular}
\end{table}

\section{S4 Rewards (all)}\label{app:s4}

\begin{table*}[t]
\caption{Fully specified S4 reward components and coefficients (Eq.~\ref{eq:s4_reward}).}
\label{tab:s4_reward}
\centering
\small
\renewcommand{\arraystretch}{1.3}
\begin{tabular}{@{}p{2.2cm}p{9.2cm}p{4.2cm}@{}}
\toprule
\textbf{Term} & \textbf{Definition / Condition} & \textbf{Coefficient(s)} \\
\midrule
$r_t^{\mathrm{progress}}$ & $w_p \cdot (d_{t-1} - d_t)$ & $w_p = 200.0$ \\[2pt]
$r_t^{\mathrm{dir}}$ & $w_f \cdot \max\!\big(0,\, \dot{\mathbf{x}}_t^{xy} \!\cdot\! \hat{v}_t\big)$, active if $d_t > 0.1\,$m & $w_f = 1.0$ \\[2pt]
$r_t^{\mathrm{heading}}$ & $w_\psi \cdot \cos(\epsilon_t)\cdot\min(\|\dot{\mathbf{x}}_t^{xy}\|, 2.0)$, active if $d_t > 0.1\,$m and $\|\dot{\mathbf{x}}_t^{xy}\| > 0.1\,$m/s & $w_\psi = 2.0$ \\[2pt]
$r_t^{\mathrm{balance}}$ & $w_b \cos\!\left(\sqrt{\phi_t^2 + \theta_t^2}\right)$ (Eq.~\ref{eq:s4_balance_speed}) & $w_b = 0.5$ \\[2pt]
$r_t^{\mathrm{speed}}$ & $w_v \exp\!\left(-\tfrac{1}{2}\left(\tfrac{\|\dot{\mathbf{x}}_t^{xy}\| - v^\star}{\sigma_v}\right)^2\right)$ (Eq.~\ref{eq:s4_balance_speed}) & $w_v = 0.2$, $v^\star = 1.0\,$m/s, $\sigma_v = 0.5\,$m/s \\[2pt]
$r_t^{\mathrm{height}}$ & $w_z \cdot \max(0,\, z_t - z_{t-1})$ & $w_z = 2.0$ \\[2pt]
$r_t^{\mathrm{bonus}}$ & $b_{\mathrm{wp}}\cdot\mathbb{1}[d_t<\tau] + b_{\mathrm{circ}}\cdot\mathbb{1}[\text{all waypoints reached}]$; circuit bonus awarded once & $b_{\mathrm{wp}} = 150.0$, $\tau = 1.5\,$m; $b_{\mathrm{circ}} = 500.0$ \\[2pt]
$r_t^{\mathrm{healthy}}$ & $w_{\mathrm{hlth}} \cdot \mathbb{1}[1.0 < z_{\mathrm{rel},t} < 2.0]$ with $z_{\mathrm{rel},t}$ from Eq.~\ref{eq:terrain_rel_z} & $w_{\mathrm{hlth}} = 5.0$ \\[2pt]
$c_t^{\mathrm{ctrl}}$ & $w_c \cdot \lVert a_t \rVert_2^2$ & $w_c = 0.1$ \\[2pt]
$c_t^{\mathrm{contact}}$ & $w_{\mathrm{cf}} \cdot \lVert F_{\mathrm{ext}} \rVert_2^2$ & $w_{\mathrm{cf}} = 5 \!\times\! 10^{-7}$ \\
\bottomrule
\end{tabular}
\vspace{2pt}
\begin{minipage}{0.98\textwidth}
\footnotesize \textit{Note:} Coefficients in this table define the raw environment reward in Eq.~\ref{eq:s4_reward}. During training, Stable-Baselines3 applies \texttt{VecNormalize} with \texttt{norm\_reward=True}, so PPO optimizes a reward scaled by running return statistics; thus nominal coefficients are not directly comparable across runs without the corresponding normalization state. During evaluation, \texttt{norm\_reward=False}, so reported episode returns are in raw reward units.
\end{minipage}
\end{table*}


\end{document}
